\section{Discussion}
\label{sec:discussion}

\paragraph{Shared inputs, model-specific mechanisms.} The most robust cross-model result is about inputs: models with
different tokenizations, depths and training objectives single out the same series for a concept, on two corpora and on
raw real windows. Whole-layer geometry is shared as well, most strongly at early, input-encoding layers, once windows are
compared without averaging away their internal order. Yet at the grain of individual causal effects concepts are
organized by model, and agreement on the same series is detected only between two closely related Chronos models. One
consistent reading is that the panel converges on \emph{what} in a series matters (its level, its seasonality) while
implementing the response differently. This is a weaker form of convergence than the representational convergence argued
for in vision and language \citep{li2016convergent,huh2024platonic}, and our instruments locate the boundary: input
agreement is measurable and replicates, whereas effect agreement across architecture families is below what ablation can
resolve at these effect sizes (Section~\ref{sec:knownanswer}).

\paragraph{What the causal evidence supports.} Single SAE features clearly beat isotropic random directions of the same
size but only modestly beat directions drawn from the activation distribution, and which features clear depends on the
intervention. The forecast-relevant information in a layer is spread over many directions, and removing any one typical
direction moves the forecast by a comparable amount. Claims that a single feature ``is'' a forecasting concept should
therefore be stated relative to an on-manifold null; against an off-manifold null a typical direction already looks
special, the time-series counterpart of the illusions documented for subspace interventions \citep{makelov2024illusion}.

\paragraph{Implications for method.} Three lessons generalize. First, activation-based selection is not a substitute for
intervention: most prominent features have no causal effect, and attention mass does not predict head importance. Second,
nulls must be validated on a known answer: a plausible uniform null inflated the count of causal features 2.7-fold.
Third, a negative verdict needs a test of sensitivity: our agreement rung never produces a false ``same'' but rarely
detects a planted shared concept, so its silence carries no information. Exploratory pipelines with thousands of
comparisons \citep{gelman2014garden} benefit from the register-then-confirm discipline used here, which confirmed between
half and two thirds of the registered claims (43 of 82 and 41 of 63). For practitioners, internal features add
information about forecast failure only for the two least accurate models; the more immediate value of the system is
diagnostic, from calibration and missing-value checks to instruments that refuse rather than return confident, wrong
results.

\section{Limitations}
\label{sec:limitations}

\paragraph{Models and alignment.} The study covers seven transformer models with contiguous tokens; models whose tokens
read disjoint lags, such as Lag-Llama, receive behavioral analysis only. With seven models, contrasts between models are
confounded with architecture, size and training data, and Timer's features do not exceed chance, so its results carry no
evidence about causal concepts. The analysis is univariate: Chronos-2's cross-series group attention is not measured
(\fid{PM-13}), and Chronos-T5's decoder and Sundial's flow head are not captured. A common time grid is no finer than the
coarsest token, and peak layer locations depend on the window size (\fid{SH-23}). Dictionaries are trained on window
vectors, which is conservative for 1-step models: trained on its own tokens, Chronos-T5-Base's dictionary explains token
variance far better and yields more causal features (\fid{MP-11}).

\paragraph{Ground truth and external validity.} Ground truth is exact only for the synthetic tier and encodes the
generators' assumptions; 12 of 33 concept parts are dominated by one real-derived generator, and overlap with the models'
pretraining data cannot be verified. Raw real benchmarks are a weaker basis for causal questions: real rows give TimesFM
representations of lower effective dimension (1.76 against 2.66) and more dead dictionary atoms (\fid{MN-20}), and weakly
predictable real-derived series receive flat forecasts in 67--100\% of cases (\fid{BM-06}; cf.\ \citealp{wan2026collapse}),
leaving an ablation nothing to move. Real data therefore enter as real-derived series (410 of 965 v2 dev series), as
dictionary-training rows and as an external check: the v3 transfer claims confirm on 800 raw real windows held out from
every dev stage.

\paragraph{Scope of the confirmation.} Held-out confirmation covers accuracy, geometry, perturbation agreement, transfer,
single-feature causal effects, atlas structure and failure prediction on two corpora. The confirmed v2 transfers derive
from two source concepts and the v3 transfers from three source models, so the transfer result rests on few concepts.
Registered claims are selected from larger dev pools, which is legitimate because the private split referees them, but
unregistered dev results remain exploratory. Per-series agreement of causal effects across architectures is not
measurable by ablation at these effect sizes, so \tool{} compares effects across models at the level of concept
profiles.

\paragraph{Dictionaries and usability.} Concept results depend on dictionary size, sparsity, revival and seed; the system
measures this dependence but cannot remove it, and feature death tracks layer geometry (alive atoms scale with effective
dimension with a log-log exponent of 1.02; \fid{MP-05}). The controls catch only failures for which a known-answer case
exists. Readability for newcomers is a design goal; no user study has measured comprehension.

\section{Conclusion}
\label{sec:conclusion}

Interpreting TSFMs is harder than it looks: tokens mean different things in different models, real series have no
annotated ground truth, and causal effects are small enough that measurement artifacts dominate unless they are
controlled. \tool{} addresses these problems with a controlled benchmark, measured time alignment, null- and reach-gated
causal concepts validated on a planted forecaster, and a register-then-confirm protocol, behind one command and a report
written for newcomers. Across seven TSFMs and two corpora, different architectures share above-null representational geometry, a
similarity that the common window average had hidden, and select the same inputs for a concept, while their causal
machinery is organized by model; causal features act mainly on level and seasonality but are only
weakly specific against on-manifold directions; activation prominence does not reveal causal importance; and
depth-resolved reaction profiles replicate for most corruptions on the primary corpus and about half on the second. Each conclusion is stated with its evidence class, and each held-out
claim was registered before the private data were opened.

\paragraph{Availability.} Code (MIT license), configurations, the example run and the findings ledger are at
\url{\repourl}. Version 1.1.0, which this paper describes, is archived
at \href{https://doi.org/10.5281/zenodo.23267462}{doi:10.5281/zenodo.23267462}, and both packages install from PyPI
(\texttt{pip install tsfm-lens tsfm-benchmark}).
